\clearpage\NoSourceHead
\begin{center}\large Zur digitalen Ausgabe\end{center}
Diese Ausgabe enthält Paul Bachmanns \textit{Über Gauss’ zahlentheoretische Arbeiten} mit Einleitung, den Abschnitten 1—27 und der Schlussbemerkung, M. A. Sterns Anhang \textit{Über einen Satz von Gauss} und die anschliessenden Bemerkungen von Bachmann und L. Schlesinger. Sie enthält ferner Ludwig Schlesingers \textit{Über Gauss’ Arbeiten zur Funktionentheorie} mit Titel, Einleitung, den zehn Hauptabschnitten, Sach- und Namenregister, historischen Berichtigungen und Inhaltsübersicht sowie den Umschlagtitel von 1933. Die Vorlage steht in \textit{Carl Friedrich Gauss, Werke}, Band X, zweite Abteilung, herausgegeben von der Gesellschaft der Wissenschaften zu Göttingen. Erfasst sind die Digitalisatblätter \SourceScope. Von Alexander Ostrowskis \textit{Über den ersten und vierten Gaussschen Beweis des Fundamentalsatzes der Algebra} sind der Titel, die Einleitung, die Abschnitte 1—6 und die historischen Berichtigungen enthalten. Das anschliessende historische Verlagsverzeichnis ist ebenfalls enthalten. Von Paul Stäckels \textit{Gauss als Geometer} sind der Umschlagtitel von 1923, der Innentitel, die Einleitung, die vollständigen Abschnitte 2 bis 34, Schlesingers Schlussbemerkung und das Inhaltsverzeichnis enthalten. Die historische Titelnotiz nennt den Abdruck aus Heft 5 der von F. Klein, M. Brendel und L. Schlesinger gesammelten Materialien und die Göttinger Nachrichten von 1917, vorgelegt am 26. Oktober 1917. Ein weiteres historisches Verlagsverzeichnis sowie Oskar Bolzas vollständige Abhandlung \textit{Gauss und die Variationsrechnung} mit Titel, Einleitung, allen vier Teilen, Schlussbemerkung und Inhaltsangabe sind ebenfalls enthalten. Philipp Maennchens Abhandlung \textit{Gauss als Zahlenrechner} ist vollständig mit historischem Umschlag, Titel, Einleitendem, den Abschnitten I—XI, Anhang und Inhaltsübersicht enthalten. Ein weiteres historisches Verlagsverzeichnis sowie Umschlag, Titel, alle fünf Teile und die Inhaltsübersicht von Harald Gepperts vollständiger Abhandlung \textit{Über Gauss’ Arbeiten zur Mechanik und Potentialtheorie} sind ebenfalls enthalten. Die abschliessenden Bemerkungen, das Gesamtinhaltsverzeichnis und das letzte historische Verlagsverzeichnis sind ebenfalls enthalten. Das Vorblatt und der Umschlag von 1922 sind enthalten; das allgemeine Titelblatt mit dem Bildnis ist noch nicht enthalten.

Die historische Titelnotiz nennt den durchgesehenen Abdruck aus Heft I der von F. Klein und M. Brendel gesammelten \textit{Materialien für eine wissenschaftliche Biographie von Gauss}, die Veröffentlichung in den \textit{Nachrichten der K. Gesellschaft der Wissenschaften zu Göttingen}, mathematisch-physikalische Klasse, 1911, und die Vorlage in der Sitzung vom 29. Juli 1911. Schlesingers Titelnotiz nennt die Neubearbeitung aus Heft III derselben Materialien, 1912, vorgelegt am 20. Juli 1912. Ostrowskis Titelnotiz nennt die Neubearbeitung des 1920 als Anhang zu Heft VIII der Materialien erschienenen Aufsatzes und das Beiheft von 1920 der Göttinger Nachrichten. Die historische Autorschaft, Zitate und bibliographischen Nachweise bleiben erhalten.

Text, Formeln und Figuren wurden von einer KI in bearbeitbares LaTeX transkribiert und gesetzt. Grundlage ist das \href{https://gdz.sub.uni-goettingen.de/id/PPN236019856}{Göttinger Digitalisat PPN236019856}; hinzugezogen wurde das \href{https://lbezone.hkust.edu.hk/obj/4/o/b761099/v10,2/pg00009.jpg}{Digitalisat der HKUST Library}. Die moderne Schrift und der neue Seitenumbruch bilden keine historische Typenidentität ab. Druckseitenzahlen stehen im Kopf, die fortlaufende Seitenzählung dieser Ausgabe unten in der Mitte.

Die historische Schreibweise, die Formeln und die Angaben des Verfassers sind beibehalten. Das Zeichen in der Beschreibung von \(\varphi(m)\) auf Druckseite 9 und in den beiden Ungleichungen auf Druckseite 24 wird entsprechend seiner sichtbaren Anordnung mit zwei waagerechten Strichen über einem Kleinerzeichen wiedergegeben. Eine mathematische oder sachliche Berichtigung des historischen Textes ist nicht Gegenstand dieser Ausgabe. Eine menschliche oder unabhängige Prüfung wird nicht beansprucht.
